
By Theorem (7.4), if $ \E X^2$ exists then 
\[
    \E X^2 = \sum_{x\in \N}^{} x^2 p_{X} (x)
.\] 

Let's see whether $ \E X^2$ exists, that is it true that 

\[
    \sum_{x \in \N}^{} |x^2| p_{X} (x) < \infty
.\] 

\begin{align*}
    \sum_{x \in \N}^{} |x^2| p_{X} (x) &= \sum_{x\in \N}^{} x^2 p_{X} (x)\\
				      &= \sum_{x \in \N}^{} x^2 \frac{\lambda^{x} e^{-\lambda} }{x!}
				      \\
				      &= \sum_{x \in \N_{\geq 1} }^{} x^2 \frac{\lambda^{x}
				      e^{-\lambda} }{x!} \\
				      &= \lambda e^{-\lambda} \sum_{x \in \N_{\geq 1}}^{} x^2
				      \frac{\lambda^{x-1} }{x!} \\
				      &= \lambda e^{-\lambda} \sum_{x \in \N_{\geq 1}}^{} x
				      \frac{\lambda^{x-1} }{(x-1)!} \\
				      &= \lambda e^{-\lambda} \sum_{x \in \N_{\geq 1} }^{} ((x-1) +
				      1) \frac{\lambda^{x-1} }{(x-1)!} \\
				      &= \lambda e^{-\lambda} \sum_{x\in \N_{\geq 1} }^{} (x-1)
				      \frac{\lambda^{x-1} }{(x-1)!} + \lambda e^{-\lambda} \sum_{x
				      \in \N_{\geq 1} }^{} \frac{\lambda^{x-1}}{(x-1)!} \\
				      &= \lambda e^{-\lambda} \sum_{x\in \N}^{} x \frac{\lambda^{x}
				      }{x!} + \lambda e^{-\lambda} \sum_{x\in \N}^{}
				      \frac{\lambda^{x} }{x!} \\
				      &= \lambda e^{-\lambda} \left( 0 + \sum_{x\in \N \geq 1}^{}
				      x \frac{\lambda^{x} }{x!} \right) + \lambda e^{-\lambda}
				      e^{\lambda} \\
				      &= \lambda^2 e^{-\lambda} \sum_{x \in \N_{\geq 1}}^{}
				      \frac{\lambda^{x-1} }{(x-1)!} + \lambda\\
				      &= \lambda^2 e^{-\lambda} e^{\lambda} + \lambda\\
				      &= \lambda^2+ \lambda < \infty
\end{align*}
We see that the expectation value exists and at the same time we find that $ \E X^2 = \lambda^2 +
\lambda$.
